% Options for packages loaded elsewhere
\PassOptionsToPackage{unicode}{hyperref}
\PassOptionsToPackage{hyphens}{url}
%
\documentclass[
]{article}
\usepackage{amsmath,amssymb}
\usepackage{iftex}
\ifPDFTeX
  \usepackage[T1]{fontenc}
  \usepackage[utf8]{inputenc}
  \usepackage{textcomp} % provide euro and other symbols
\else % if luatex or xetex
  \usepackage{unicode-math} % this also loads fontspec
  \defaultfontfeatures{Scale=MatchLowercase}
  \defaultfontfeatures[\rmfamily]{Ligatures=TeX,Scale=1}
\fi
\usepackage{lmodern}
\ifPDFTeX\else
  % xetex/luatex font selection
\fi
% Use upquote if available, for straight quotes in verbatim environments
\IfFileExists{upquote.sty}{\usepackage{upquote}}{}
\IfFileExists{microtype.sty}{% use microtype if available
  \usepackage[]{microtype}
  \UseMicrotypeSet[protrusion]{basicmath} % disable protrusion for tt fonts
}{}
\makeatletter
\@ifundefined{KOMAClassName}{% if non-KOMA class
  \IfFileExists{parskip.sty}{%
    \usepackage{parskip}
  }{% else
    \setlength{\parindent}{0pt}
    \setlength{\parskip}{6pt plus 2pt minus 1pt}}
}{% if KOMA class
  \KOMAoptions{parskip=half}}
\makeatother
\usepackage{xcolor}
\setlength{\emergencystretch}{3em} % prevent overfull lines
\providecommand{\tightlist}{%
  \setlength{\itemsep}{0pt}\setlength{\parskip}{0pt}}
\setcounter{secnumdepth}{-\maxdimen} % remove section numbering
\ifLuaTeX
  \usepackage{selnolig}  % disable illegal ligatures
\fi
\usepackage{bookmark}
\IfFileExists{xurl.sty}{\usepackage{xurl}}{} % add URL line breaks if available
\urlstyle{same}
\hypersetup{
  hidelinks,
  pdfcreator={LaTeX via pandoc}}

\author{}
\date{}

\begin{document}

\section{Nef bounded products and both derivatives for continuous
obstacles}\label{nef-bounded-products-and-both-derivatives-for-continuous-obstacles}

This paper constructs the actual raw bounded metric products and proves
both derivatives of continuous-obstacle energy for every nef rational
line bundle in the stated domain. Independent review122 returned PASS
for this exact bounded scope. The original general P16.5 problem remains
open: no complete unbounded finite-energy equation or general nonnef
reference theorem is asserted here.

\subsection{1. Actual scope and the corrected source
boundary}\label{actual-scope-and-the-corrected-source-boundary}

Let k be algebraically closed, trivially valued and of characteristic
zero, with arbitrary cardinality. Let X be smooth, integral and
projective of dimension \(n\geq1\). Retain an actual nef rational line
bundle L, its numerical class theta, and its canonical trivial reference
zero. Fix an actual ample rational bundle H with class h. All potentials
below belong to the genuine numerical cone of BJ Definition4.1; no
quotient cone or inverse-lift operator is defined.

The primary source is Boucksom--Jonsson, arXiv:1801.08229v4, header
March2022, 143pages. Theorem5.20 itself states the envelope property for
\textbf{every nef theta}, not just ample theta, under the retained
smoothness and characteristic assumptions. Its printed proof reduces to
ample rational classes through Lemma5.9; that reduction does not narrow
the theorem's conclusion. Complete printed statement/proof units checked
here are local literal LF2980--2985 and3027--3062, and fresh native-PDF
lines4269--4274 and4321--4363. Lemma5.9's complete proof is already
checked at LF2801--2842. Published recursive inputs in the
multiplier-ideal argument are retained, not independently recursively
certified.

Two distinctions matter. Lemma5.17's equivalence with continuity of
continuous-obstacle envelopes is formulated for ample classes.
Theorem5.20 does not by that citation alone prove continuity of every
nef continuous-obstacle envelope. Remark7.16 further warns that the
floor-defined E1(theta) can strictly contain the ambient finite-energy
vector space intersected with PSH(theta). Theorems7.14/7.18 and
Proposition7.19 therefore cannot simply be applied to every member of
E1(theta).

The exact accepted ample G0 theorem is the actual-ample variational
theorem on the original projective analytic space, with canonical
reference and raw top-intersection normalization. Its exact proof
provenance and SHA256 are recorded with this paper. Only its actual
ample projective domain is a grant. Its raw energy, genuine metric
operator, energy concavity, continuous envelopes, and both
continuous-direction outer derivatives are used at theta+epsilon h. No
nonample equation or unbounded nonample continuity is granted.

\subsection{2. Bounded genuine potentials fit the actual published
product
domain}\label{bounded-genuine-potentials-fit-the-actual-published-product-domain}

Write B(theta)=PSH(theta) intersect bounded functions. This is a
statement about the entire genuine numerical bounded cone, not a newly
defined smaller test family. Since theta is nef, zero and every constant
are genuine theta-psh by BJ4.3. If u belongs to B(theta), then u belongs
to PSH(theta+h), since zero is h-psh and numerical addition is BJ4.7. It
is bounded, so its ample energy is finite by monotonicity between its
constant bounds. Thus the restriction of u to X\^{}lin belongs to the
actual ambient vector space E\textasciitilde1 of Definition7.12. The
bounded potential itself remains a function on the entire original
X\^{}an; the auxiliary restriction is only the published pairing domain.

Theorems7.14 and7.18 now apply with their stated hypotheses. Define the
raw bounded energy and mixed measures by their actual energy-pairing
formulas:

\[E_\theta(u)=\frac{(\theta,u)^{n+1}-(\theta,0)^{n+1}}{n+1},\qquad
M_\theta(u_1,\ldots,u_n)=\bigwedge_{i=1}^n(\theta+dd^cu_i).\]

The total mass is theta\^{}n, and M is positive for genuine bounded
theta-psh inputs, by Theorem7.18 and Proposition7.19(i). The energy
shifts by c theta\^{}n under u -\textgreater{} u+c.~At a finite PL/model
stage, Proposition7.19(ii) gives the actual sum of component Dirac
measures with coefficients

\[\operatorname{ord}_F(\mathcal X_0)\prod_i(\theta_{\mathcal X}+D_i)|_F.\]

These are actual model metric intersection products with the given
reference, component multiplicities and bundle class. For continuous nef
inputs, CLD compatibility follows by the actual metric decomposition
(L,u)=(L+H,u)-(H,0). Each term has an actual ample class and a verified
continuous psh metric. BJ Remark7.23 applies to those ample terms; BE v3
Sections7.14 and8.4--8.7 provide the multilinear difference-of-psh
extension and regularization. Expanding the actual metric differences
identifies the continuous nef products. Remark7.23 alone is not cited as
a nonample statement. Individual current factors are not asserted to
have been constructed on all bounded inputs merely from the wedge
notation.

One can also recover the same bounded theta-products as limits of actual
positive ample operators before taking any limit. For
epsilon\textgreater0 rational, each u\_i remains genuinely
(theta+epsilon h)-psh. The actual G0 ample operator on L+epsilon H is
defined, positive, and has raw mass (theta+epsilon
h)\^{}n.~Multilinearity of the published mixed measures makes its
pairing with each continuous test a finite polynomial in epsilon, with
the constant term M\_theta(u\_1,\ldots,u\_n). Thus the positive ample
products converge weakly to the displayed genuine theta-product.
Uniqueness in Theorem7.18 makes this independent of H. The limit is
identified with an already constructed original metric product; it is
not the definition of an arbitrary inverse lift.

For a decreasing net u\_i in B(theta) with bounded genuine limit u, any
tail has common upper and lower bounds. All terms lie in E1(theta+h),
and Proposition7.19(iv) yields the original bounded mixed weak
continuity. Theorem7.14 gives bounded-energy continuity for these nets.
Integrability and the energy-pairing identity for an ambient
E\textasciitilde1 test are Proposition7.19(v). This proves no weighted
theorem for a test outside that stated ambient domain.

\subsection{3. Genuine envelopes at the nef
boundary}\label{genuine-envelopes-at-the-nef-boundary}

For clarity we can prove the nef class-closure step directly, without a
general abstract monotone-net diagonal assertion. There is an equality
of genuine cones

\[\operatorname{PSH}(\theta)=\bigcap_{\epsilon\in\mathbb Q_{>0}}\operatorname{PSH}(\theta+\epsilon h).\]

The forward inclusion is numerical addition with the genuine constant
h-psh metric. For the converse, the actual ample approximation theorem
gives finite continuous FS functions f\_\{epsilon,j\} decreasing to a
member u of every displayed ample cone. Use triples (epsilon,j,delta),
with delta\textgreater0, and functions f\_\{epsilon,j\}+delta. Given two
triples, choose a smaller positive epsilon\_3 and delta\_3 smaller than
both previous deltas. Their old minimum b is finite continuous, and
u+delta\_3\textless b at every point. The increasing open sets
\{f\_\{epsilon\_3,j\}+delta\_3\textless b\} cover compact X\^{}an. A
finite subcover, followed by one common later net index, provides a
third triple dominated by both old functions. Including epsilon and
delta as ordered coordinates makes this a decreasing defining FS net.
Refining any prescribed finite pointwise threshold, including at a pole,
proves its endpoint is u. Its classes approach theta, so the numerical
definition gives u in PSH(theta).

Consequently a decreasing family whose classes approach theta and which
stays above a finite constant belongs, at its endpoint, to every fixed
approaching ample cone and hence to PSH(theta). This proves the
particular closure needed below. It does not claim a proof of the
source's general abstract Lemma4.6.

For continuous f on the original compact X\^{}an, let P\_theta f be the
pointwise supremum of genuine theta-psh functions bounded by f.~Constant
competitors show that it lies between min f and max f.~The nef envelope
property makes the upper regularization theta-psh; continuity of f makes
that regularization \textless=f, so it is itself a competitor. Hence
P\_theta f equals that regularization, is genuinely theta-psh, and is
bounded. This argument does not assert its continuity.

Put theta\_epsilon=theta+epsilon h and q\_epsilon=P\_theta\_epsilon
f.~The actual ample theory makes q\_epsilon continuous. As epsilon
decreases to zero, q\_epsilon decreases pointwise. The common constant
bounds make the limit generically finite, and the preceding genuine nef
intersection proves that its limit is theta-psh. That limit is
\textless=f.~Conversely every theta-psh competitor is
theta\_epsilon-psh, so P\_theta f\textless=q\_epsilon. Therefore

\[q_\epsilon\downarrow q_0=P_\theta f.\]

All q\_epsilon, for epsilon\textless=1, lie in the one actual ample cone
PSH(theta+h) and have the same two constant bounds. Theorems7.14/7.18
and Proposition7.19 therefore apply to this decreasing family within
their proper domain. Polynomial expansion, coefficient by coefficient,
proves

\[E_{\theta_\epsilon}(q_\epsilon)\longrightarrow E_\theta(q_0),\qquad
M_{\theta_\epsilon}(q_\epsilon)\longrightarrow M_\theta(q_0).\]

Some individual theta-coefficients at positive epsilon may be signed;
positivity of those coefficients is neither used nor asserted. Their
continuous-test pairings converge by Proposition7.19(iv) in the common
ample finite-energy cone. The full operators on the left and the final
operator on the right are positive. For each signed measure coefficient,
expand every factor (theta,u) as (theta+h,u)-(h,0). This is a fixed
finite sum of positive mixed products, with their signs retained. Each
positive summand has the fixed finite numerical mixed mass of its ample
classes, independent of u and epsilon. The sum of those masses bounds
total variation uniformly, so multiplying a coefficient by a positive
power of epsilon makes it tend weakly to zero. The corresponding energy
coefficients converge by the bounded published pairing theorem in the
same common ample cone. This justifies the varying-epsilon polynomial
passage with actual signed coefficients, rather than only PL test
functions.

\subsection{4. Uniform continuity needed for
variation}\label{uniform-continuity-needed-for-variation}

Pointwise order and constants give the original envelope estimate

\[\|P_\theta f-P_\theta g\|_\infty\leq\|f-g\|_\infty.\]

For uniformly converging bounded genuine theta-psh potentials u\_j
-\textgreater{} u, their theta mixed measures and energy converge. Here
is a reduction to the published decreasing-net theorem that avoids any
unproved nonample strong topology. From any sequence converging in sup
norm, pass to a subsequence with errors at most \(4^{-j}\); adding
\(3\cdot4^{-j}\) makes that subsequence decreasing pointwise to u.
Constants do not change curvature products. Energy changes only by
\(3\cdot4^{-j}\theta^n\). The decreasing convergence in the common ample
cone from Section2 therefore gives the desired convergence on that
subsequence. Every subsequence admits such a further subsequence, so the
entire sequence converges. For each fixed continuous measure test this
proves continuity on the sup-norm metric space; consequently it also
proves convergence for arbitrary sup-norm converging nets. It does not
apply a false monotone convergence theorem to arbitrary Borel nets.

\subsection{5. Both derivatives of continuous-obstacle
energy}\label{both-derivatives-of-continuous-obstacle-energy}

Define F\_theta(f)=E\_theta(P\_theta f) for f in C0(X\^{}an). It is
finite, concave and translation equivariant with raw coefficient
theta\^{}n.~Concavity follows from the genuine convex cone, envelope
order and bounded energy concavity; alternatively it follows from the
pointwise limit of the ample concave functionals in Section3.

We first give the needed ample supporting inequalities rather than
importing an unproved assertion about arbitrary obstacles. For an ample
class omega and continuous f, q=P\_omega f is continuous genuine. If
0\textless=s\textless=1, monotonicity gives P\_omega(q+s(f-q))=q. The
actual ample G0 outer derivative at q in the continuous direction f-q
therefore yields orthogonality:

\[\int(f-q)M_\omega(q)=0.\]

Energy concavity along genuine psh segments gives

\[E_\omega(r)\leq E_\omega(q)+\int(r-q)M_\omega(q),\qquad r=P_\omega g.\]
Since r\textless=g and M\_omega(q) is positive, the displayed
orthogonality proves the supporting inequality

\[F_\omega(g)-F_\omega(f)\leq\int(g-f)M_\omega(P_\omega f).\]

Use this inequality once at f and once at f+tg. For t\textgreater0 it
bounds the quotient between integral g M\_omega(P\_omega(f+tg)) and
integral g M\_omega(P\_omega f). Apply it with omega=theta\_epsilon and
pass epsilon down to zero using the actual energy and operator limits
from Section3. Thus

\[\int gM_\theta(P_\theta(f+tg))\leq
\frac{F_\theta(f+tg)-F_\theta(f)}t\leq
\int gM_\theta(P_\theta f).\]

The envelope Lipschitz estimate and Section4 make the left integral
converge to the right one as t down to zero. For t\textless0 the same
two supporting inequalities reverse under division; the two bounds have
the same limit as t increases to zero. Therefore both derivatives exist
and agree:

\[\left.\frac{d}{dt}\right|_{t=0}F_\theta(f+tg)=\int gM_\theta(P_\theta f),\qquad f,g\in C^0(X^{\rm an}).\]

In particular this proves both outer signs at every \textbf{continuous
genuine} theta-psh base potential u, by taking f=u. It proves the
continuous-obstacle formula for every nef actual L, without
semiampleness and including zero raw top mass. It does not infer a
derivative at an arbitrary unbounded member of the floor-defined
E1(theta), or even at an arbitrary bounded discontinuous base, from this
continuous-obstacle argument. For example, on the product of two
projective lines with L pulled back from one factor, the raw top mass is
zero and the present top-degree energy is constant. A different
lower-degree mixed energy can have positive mass and is a different
functional.

\subsection{6. Remaining original
obligations}\label{remaining-original-obligations}

This supplies a genuine bounded top-degree operator, nef envelope
membership, and both derivatives for continuous obstacles beyond
semiample contractions. It does not supply the complete unbounded
E1(theta) operator with all original weighted/floor limits, intrinsic
target-space characterization, or variational attainment and uniqueness
for every finite-energy measure. The nef raw top mass is theta\^{}n;
proving all required section-volume and reference identifications in the
intended full general class remains part of the original problem.
General nonnef big references and general nonbig nonzero classes are
also still unresolved here.

The explicit published source domains are the important constraint.
Remark7.16 prohibits silently identifying the floor-defined E1(theta)
with E\textasciitilde1 intersect PSH(theta). A bounded approximation
theorem cannot be promoted to that larger class without an energy
estimate controlling the full limiting operator. These bounded
conclusions received the scoped independent PASS identified below.
Original P16.5 remains open.

\subsection{7. Source/application
ledger}\label{sourceapplication-ledger}

\begin{itemize}
\tightlist
\item
  BJv4 LF2190--2237: genuine numerical definition, joint decreasing
  class closure, constants for nef theta. LF2330--2368: addition and
  uniform closure proof. Used for Sections2--4.
\item
  BJv4 LF2801--2842: Lemma5.9 and Theorem5.11 complete
  statements/proofs. LF2980--2985 and3027--3062: complete Theorem5.20
  statement and proof. Fresh native-PDF selected units are retained as
  source audit only. Used for nef envelope membership, not unproved nef
  continuity.
\item
  BJv4 LF2909--2935: complete Lemma5.17 statement/proof, whose ample
  domain is retained. Used only for theta\_epsilon continuous envelopes.
\item
  BJv4 LF3890--3955: actual ambient E\textasciitilde1 definition,
  complete Theorem7.14 and Lemma7.15 proofs, and Remark7.16.
  LF3980--4057: complete Theorem7.18 and Proposition7.19
  statements/proofs. Used only on verified bounded inputs in a common
  ample finite-energy cone.
\item
  Exact whole ample G0 body named in Section1: actual raw operator,
  energy concavity and both outer derivatives within its accepted ample
  domain. Every use at positive epsilon retains that domain. No original
  P16.5 or other unproved nonample premise is a grant.
\item
  Primary URL: https://arxiv.org/pdf/1801.08229v4 . Whole-book and
  recursive-foundation verification are not claimed. No finite numerical
  check certifies the deduction.
\end{itemize}

Independent review122 checked 17 external primary premises with 29 HTTPS
source passages, and one separately indexed inline task grant for the
exact actual-ample G0 theorem. The latter passage uses the task URI
urn:xean:task and is not an external web citation. All 30 literal
passage statements match their indexed premises. The canonical report is
preserved without replacing its conclusion or inventing additional
source passages.

The source-domain clarifications incorporated here retain the X\^{}lin
restriction only for the ambient pairing membership and retain the ample
domain of Remark7.23. The nef continuous compatibility uses the explicit
actual metric difference above. No continuity of all nef envelopes, or
equality of the larger intrinsic E1(theta) with ambient
E\textasciitilde1, is inferred. All-cardinality decreasing nets and the
full original continuous test space are retained.

\end{document}
